\documentclass[11pt]{article}

% ============================================================================
%  Rally-Racing — System Guide (LaTeX source)
%
%  This is the editable source for the HTML system guide opened by the GUI's
%  "System Guide" button. To regenerate the HTML after editing:
%
%      make4ht system_guide.tex            (TeX4ht; produces system_guide.html)
%  or  pandoc system_guide.tex -s -o system_guide.html --mathjax
%
%  The repository ships a pre-built docs/system_guide.html so end users do not
%  need a LaTeX toolchain. Rebuild only if you change this source.
% ============================================================================

\usepackage[utf8]{inputenc}
\usepackage{amsmath}
\usepackage{amssymb}
\usepackage{hyperref}
\usepackage{geometry}
\geometry{margin=1in}

\title{Rally-Racing --- System Guide}
\author{Operation, Architecture, and Function Reference}
\date{}

\begin{document}
\maketitle

\tableofcontents

% ----------------------------------------------------------------------------
\section{Overview}

Rally-Racing is a PyBullet driving environment trained with Proximal Policy
Optimization (PPO). A car learns to pass a sequence of checkpoints while
avoiding obstacles and handling ramps. The project supports two observation
regimes:

\begin{itemize}
  \item \textbf{Privileged observation} --- the policy receives the
        ground-truth position of the nearest obstacle directly from the
        simulator.
  \item \textbf{Vision observation} --- a convolutional network
        (\texttt{ObstacleCNN}) infers obstacle presence and location from the
        car's camera, and those predictions replace the privileged obstacle
        channels. A privileged policy is fine-tuned into a vision policy.
\end{itemize}

Training follows a curriculum of increasingly difficult scenarios, and the
whole system can be operated either from the command line (manual operation)
or through the Tkinter control panel (GUI operation). Both are described in
Section~\ref{sec:operation}.

% ----------------------------------------------------------------------------
\section{The Curriculum}

The scenarios build on one another. Each privileged policy warm-starts from
the previous scenario's privileged policy.

\begin{center}
\begin{tabular}{ll}
\textbf{Scenario} & \textbf{Track contents} \\
\hline
phase1 & 6 checkpoints, no obstacles, no ramps \\
phase2 & 6 checkpoints + 4 obstacles \\
phase3 & 6 checkpoints + 4 obstacles + 2 ramps \\
custom & bespoke checkpoint / obstacle / ramp lists \\
\end{tabular}
\end{center}

Warm-start chain: $\text{phase1} \rightarrow \text{phase2} \rightarrow
\text{phase3} \rightarrow \text{custom}$. A \emph{vision} policy for a given
scenario warm-starts from that scenario's \emph{privileged} policy, not from a
previous vision policy.

% ----------------------------------------------------------------------------
\section{Operation}\label{sec:operation}

\subsection{GUI operation (recommended)}

The control panel (\texttt{src/control.py}) exposes every setting and runs
every pipeline stage without touching source code. Launch it by
double-clicking the \emph{Rally-Racing Control Panel} icon (after running
\texttt{install\_desktop.sh} once), or from a terminal:

\begin{verbatim}
./launch.sh
\end{verbatim}

The panel is organised into a settings area (left) and a run area (right):

\begin{itemize}
  \item \textbf{Training tab} --- total timesteps, parallel environments,
        scenario, PPO hyperparameters, and Weights \& Biases logging.
  \item \textbf{Reward weights tab} --- every coefficient in
        \texttt{RewardConfig} (Section~\ref{sec:reward}), editable live.
  \item \textbf{Vision tab} --- dataset collection size, CNN training
        settings, and vision fine-tune budget.
  \item \textbf{Evaluation tab} --- scenarios to evaluate, episode count,
        whether to open the simulator window, and the model file to load.
  \item \textbf{Run \& Console tab} --- run any stage individually, or tick
        several and run them in succession; live output streams to the
        console.
  \item \textbf{Training Metrics tab} --- live plots of episode reward, loss,
        and other scalars read from the local TensorBoard logs, plus an
        \emph{Open in W\&B} button for the full online dashboard.
\end{itemize}

Settings are saved to \texttt{gui\_config.json}, which the training scripts
read at runtime. Nothing is hardcoded in two places.

\subsection{Manual operation (command line)}

Every stage is also a standalone CLI script. The GUI is a convenience layer
over exactly these commands; the scripts remain fully usable on their own.

\paragraph{Privileged training.}
\begin{verbatim}
python3 src/train.py --scenario phase1
python3 src/train.py --scenario phase3 --timesteps 150000 --no-wandb
python3 src/train.py --scenario phase3 --fresh   # ignore warm-start chain
\end{verbatim}

\paragraph{Vision fine-tune} (requires the matching privileged model):
\begin{verbatim}
python3 src/train_vision.py --scenario phase3
python3 src/train_vision.py --scenario phase3 --timesteps 200000
\end{verbatim}

\paragraph{Evaluation} (omit \texttt{--no-render} to watch in a window):
\begin{verbatim}
python3 src/test.py --model gen2_models/phase3/privileged/best/best_model.zip \
    --scenarios phase3 --episodes 20 --no-render
python3 src/test_vision.py --ppo gen2_models/phase3/vision/best/best_model.zip \
    --scenarios phase3 --episodes 20 --no-render
\end{verbatim}

\paragraph{Relationship between the two.} A GUI ``Run'' button builds and
launches the corresponding command above as a subprocess, passing your saved
settings via \texttt{gui\_config.json} and command-line arguments. Manual and
GUI operation therefore produce identical runs; the GUI simply removes the
need to remember flags and paths.

% ----------------------------------------------------------------------------
\section{Architecture and Key Functions}

This section documents the important functions across the package, grouped by
module.

\subsection{Environment: \texttt{simple\_driving/envs/rally\_driving\_env.py}}

\paragraph{\texttt{RallyDrivingEnv.reset(seed, options)}}
Re-initialises an episode: respawns the car at the scenario's start pose,
rebuilds the checkpoint list, and (for phase2/phase3) spawns obstacles and
ramps. Returns the initial observation. The random seed makes obstacle
placement reproducible.

\paragraph{\texttt{RallyDrivingEnv.step(action)}}
Advances the simulation by one control step. It applies the action, steps
PyBullet, computes the reward through the injected
\texttt{reward\_callback}, checks termination, and returns the tuple
$(\text{obs}, \text{reward}, \text{done}, \text{truncated}, \text{info})$.
If no reward callback was supplied it raises, by design --- the environment
never silently scores episodes with a default reward.

\paragraph{\texttt{RallyDrivingEnv.getExtendedObservation()}}
Builds the 11-dimensional observation vector the policy consumes:
\[
[\, g_x,\ g_y,\ c_{\text{rem}},\ o_x,\ o_y,\ b_{\text{obs}},\
      \theta_{\text{pitch}},\ \theta_{\text{roll}},\ v,\
      \cos\psi,\ \sin\psi \,]
\]
where $(g_x,g_y)$ is the next checkpoint in the car's local frame,
$c_{\text{rem}}$ is the fraction of checkpoints remaining, $(o_x,o_y)$ is the
nearest obstacle in the local frame, $b_{\text{obs}}$ is a $0/1$ flag for
whether an obstacle exists, $\theta_{\text{pitch}}$ and $\theta_{\text{roll}}$
are chassis angles, $v$ is planar speed, and $(\cos\psi,\sin\psi)$ encodes
heading without a wrap discontinuity. Transforming the goal and obstacle into
the car's local frame makes the policy translation- and rotation-invariant.

\paragraph{\texttt{RallyDrivingEnv.get\_camera\_image(width, height)}}
Renders the car's onboard camera view. Purely additive: it does not affect
\texttt{step}, \texttt{reset}, or the reward, so it can be used for vision data
collection without altering training dynamics.

\paragraph{\texttt{RallyDrivingEnv.\_nearest\_obstacle(car\_pos)}}
Returns the closest obstacle and its distance, used both for the observation
and for the close-contact reward term.

\subsection{Reward: \texttt{src/reward.py}}\label{sec:reward}

\paragraph{\texttt{RewardConfig}}
A dataclass-like container of every reward coefficient (goal reward, obstacle
penalty, progress scale, smoothness penalties, repulsive-field parameters,
airborne bonus). The GUI edits these values; the environment never hardcodes
them.

\paragraph{\texttt{RewardCalculator.\_\_call\_\_(\dots)}}
Computes the scalar reward for one step. The total is a sum of independent
terms:
\[
r = r_{\text{prog}} + r_{\text{reg}} + r_{\text{goal}} + r_{\text{smooth}}
  + r_{\text{coll}} + r_{\text{rep}} + r_{\text{oob}} + r_{\text{air}}
\]
where each term contributes as follows:
\begin{itemize}
  \item \textbf{Progress} $r_{\text{prog}} = r_{\text{step}} + k_{\text{prog}}\,\Delta d$,
        a constant per-step cost plus a bonus proportional to the reduction in
        distance to the goal $\Delta d$.
  \item \textbf{Regression} $r_{\text{reg}} = r_{\text{reg}}$ when $\Delta d < 0$,
        an extra penalty for moving away from the goal. It fires in all
        scenarios (including phase1) so the agent always pays for zigzagging.
  \item \textbf{Checkpoint} $r_{\text{goal}}$ added once when a checkpoint is
        reached.
  \item \textbf{Smoothness / stability}
        $r_{\text{smooth}} = k_{\text{yaw}}\,|j_\psi|
         + k_{\text{roll}}\,|\Delta\theta_r| + k_{\text{pitch}}\,|\Delta\theta_p|$,
        penalising yaw jerk $j_\psi$ (the change in yaw rate) and roll / pitch
        changes between steps.
  \item \textbf{Collision} $r_{\text{coll}}$ applied when the nearest obstacle
        is within the minimum safe distance.
  \item \textbf{Repulsive field}
        $r_{\text{rep}} = -k_{\text{rep}} \sum_i \max(0,\ R - d_i)$,
        a soft penalty that grows as the car approaches any obstacle within
        radius $R$, discouraging skimming even without a collision.
  \item \textbf{Boundary} $r_{\text{oob}}$ applied when the car leaves the
        world bounds.
  \item \textbf{Airborne} $r_{\text{air}}$, a small phase3 bonus for being
        pitched up while still making forward progress (rewarding committed
        ramp jumps).
\end{itemize}

\paragraph{\texttt{RewardCalculator.\_wrap\_delta(delta)}}
Wraps an angular difference into $[-\pi, \pi]$ so heading changes near the
$\pm\pi$ boundary are measured correctly rather than as near-full rotations.

\paragraph{\texttt{custom\_reward(**kwargs)}}
Module-level functional wrapper around a default \texttt{RewardCalculator},
used as the environment's \texttt{reward\_callback}.

\subsection{Vision: \texttt{vision/model.py}}

\paragraph{\texttt{ObstacleCNN.forward(x)}}
Runs the convolutional stack over a camera frame and returns the raw network
outputs (logits / regressed coordinates) describing whether an obstacle is
visible and where.

\paragraph{\texttt{ObstacleCNN.predict(x)}}
Inference wrapper around \texttt{forward} that converts raw outputs into the
obstacle-presence flag and local position estimate the vision environment
substitutes into the observation vector.

\subsection{Vision environment: \texttt{simple\_driving/envs/vision\_rally\_env.py}}

\paragraph{\texttt{VisionRallyDrivingEnv.getExtendedObservation()}}
Overrides the privileged observation: it captures a camera frame, runs the
CNN, and replaces the obstacle channels the obstacle channels $(o_x, o_y, b_{\text{obs}})$
with the network's prediction while leaving all other channels identical. This
is what lets a privileged-trained policy be fine-tuned to drive from vision
with minimal architectural change.

\subsection{Training: \texttt{src/train.py}}

\paragraph{\texttt{warm\_start\_path(scenario)}}
Returns the path of the policy a new run should warm-start from, encoding the
curriculum chain. Returns \texttt{None} when training should start from
scratch (e.g. phase1).

\paragraph{\texttt{make\_train\_env(scenario, n\_envs)} /
\texttt{make\_eval\_env(scenario, log\_dir)}}
Construct the vectorised training environments and the monitored evaluation
environment, respectively, wiring in the reward callback and scenario
configuration.

\paragraph{\texttt{make\_callbacks(eval\_env, best\_dir, log\_dir, use\_wandb, n\_envs)}}
Assembles the Stable-Baselines3 callbacks: periodic evaluation with best-model
saving, and (optionally) the Weights \& Biases callback. The
best-by-evaluation snapshot is what evaluation later loads.

\paragraph{\texttt{main()}}
Orchestrates a training run: parse arguments, build environments, load or
create the PPO model, attach callbacks and logging, train for the requested
timesteps, and save the final model. When W\&B is enabled it also records the
run URL to \texttt{logs/<run>/wandb\_url.txt} so the GUI can open the exact run.

\subsection{Evaluation: \texttt{src/test.py}}

\paragraph{\texttt{evaluate(model\_path, scenarios, render, episodes)}}
Loads a saved policy and runs it for the requested number of episodes per
scenario, reporting per-episode reward, step count, and checkpoint completion.
With \texttt{render=True} it opens a PyBullet window so the car can be watched;
with \texttt{render=False} it runs headless for fast batch evaluation.

% ----------------------------------------------------------------------------
\section{Logging and Metrics}

Training writes TensorBoard scalar logs to
\texttt{logs/<gen>\_<scenario>\_<observation>/}. Because training initialises
W\&B with \texttt{sync\_tensorboard=True}, those same scalars also sync to the
configured W\&B project when logging is enabled. The GUI's
\emph{Training Metrics} tab reads the local event files directly (no W\&B login
required) and plots them live; the \emph{Open in W\&B} button jumps to the
online run when a URL was recorded.

% ----------------------------------------------------------------------------
\section{Known Behaviour}

\begin{itemize}
  \item \textbf{Phase3 privileged policy is roughly 90\% reliable.} Over many
        episodes most are clean laps, with occasional catastrophic failures in
        awkward ramp or obstacle configurations.
  \item \textbf{Phase1 evaluation is deterministic.} With no obstacles and a
        fixed spawn, a deterministic policy produces identical episodes ---
        identical eval rewards are expected, not a bug.
  \item \textbf{Vision policies score somewhat below privileged.} The
        CNN-derived observation is noisier than ground truth. A large gap
        suggests the CNN is not engaging; check the \texttt{vision\_active}
        percentage reported by \texttt{test\_vision.py}.
\end{itemize}

\end{document}
